% =============================================================================
% CHAPTER 3: Cosmic Chronology — From Inflation to the Current Beat
% Sources: chronology.md, theory_v4_complete.md §4, post cosmic-chronology
% =============================================================================

\chapter{Cosmic Chronology: From Inflation to the Current Beat}
\label{ch:chronology}

The evolution of brane tension from the Big Bang to today reveals how the universe tuned itself to its fundamental frequency.

% -----------------------------------------------------------------------------
\section{Timeline of Brane Evolution}
\label{sec:timeline}

\begin{table}[h]
\centering
\caption{Evolution of brane tension through cosmic history.}
\label{tab:timeline}
\begin{tabular}{@{}llll@{}}
\toprule
\textbf{Phase} & \textbf{Age} & \textbf{$\tau$ (J/m$^2$)} & \textbf{Description} \\
\midrule
Inflation & $0 \to 10^{-34}$~s & $10^{50}$ & Quasi-exponential expansion, hyper-tense brane \\
Brane Reheating & $10^{-34} \to 10^{-32}$~s & $10^{30}$ & Tension decay via MN-anti-MN production in bulk \\
Relaxation & $10^{-32}$~s $\to$ 1~Gyr & $10^{27} \to 7 \times 10^{19}$ & $\tau \propto t^{-1/2}$, fundamental mode enters resonance \\
Current Era & 13.8~Gyr & $7 \times 10^{19}$ & Stable oscillation with $2\Gyr$ period \\
\bottomrule
\end{tabular}
\end{table}

% -----------------------------------------------------------------------------
\section{The Violent Birth}
\label{sec:violent_birth}

In this framework, the brane appears at the Big Bang with quasi-Planckian tension $\tau_\text{BB} \sim 10^{50}$~J/m$^2$---a membrane stretched to breaking point, vibrating with pure energy.

\paragraph{Phase I --- Trans-membrane Inflation ($0$--$10^{-34}$~s).}
The colossal excess tension fuels exponential expansion. The membrane expands like a soap bubble blown by a hurricane, creating space from dimensional nothingness. The extreme curvature prevents any oscillatory modes.

\paragraph{Phase II --- Brane Reheating ($10^{-34}$--$10^{-32}$~s).}
As inflation ends, the brane tension converts to particle production. Massive MN-anti-MN pairs are created in the bulk. Energy density transfers from the geometric to the matter sector. Tension drops abruptly by 20 orders of magnitude, leaving residual tension around $10^{30}$~J/m$^2$.

\paragraph{Phase III --- Slow Stabilization ($10^{-32}$~s--$100$~Myr).}
Tension relaxes logarithmically toward its current value. Like a violin string being tuned, the membrane seeks its natural frequency. The brane tension follows a power-law decay:
\begin{equation}
\tau(t) = \tauZ \left(\frac{t_0}{t}\right)^{1/2}
\end{equation}

This natural cooling allows the fundamental mode to enter resonance when the oscillation period matches the age of the universe.

% -----------------------------------------------------------------------------
\section{The Awakening of Oscillations}
\label{sec:awakening}

Only when $\tau$ becomes ``loose enough'' does the fundamental mode enter the $T \sim 2\Gyr$ band. Oscillation starts about 1~Gyr after the Big Bang---exactly when DESI's baryon acoustic oscillations and Planck's ISW resonance independently confirm the fundamental period.

This temporal coincidence is no accident: it is the moment when the universe, finally tuned, begins playing its fundamental melody. During this critical epoch:

\begin{itemize}
  \item The conformal symmetry ($T^\mu{}_\mu = 0$) of the radiation-dominated era has already been broken by the QCD phase transition at $\sim 257\MeV$.
  \item The trace coupling factor $(1 - 3\weff)$ has switched from 0 to 1.
  \item The Cosmic Web has begun to form, providing the macroscopic forcing $\Fweb[E_{\mu\nu}]$.
  \item The ER=EPR-entangled PBH network is already in place from primordial formation.
\end{itemize}

Today, the brane has reached its equilibrium configuration: stable tension $\tauZ = 7 \times 10^{19}$~J/m$^2$, fundamental period $T = 2.0\Gyr$, with $\sim 10\%$ of dark matter participating in oscillations.

% -----------------------------------------------------------------------------
\section{Tension Calibration: The Perfect Tuning}
\label{sec:calibration}

The period is not given by a simple harmonic formula. The stick-slip cycle has:
\begin{equation}
T \approx t_\text{stick} + t_\text{slip} \approx 2.0\Gyr
\end{equation}
where $t_\text{stick}$ is the charging time ($E_{\mu\nu}$ forcing against the Goldberger--Wise restoring potential) and $t_\text{slip}$ is the rapid discharge time. The harmonic approximation:
\begin{equation}
T \approx 2\pi\sqrt{\frac{f_\text{osc}\,M_\text{DM,tot}}{\tauZ}}
\end{equation}
gives the correct order of magnitude, but the precise period requires numerical integration of the full V8.0 ODE including the $\xi R\phi$ attractor term.

Inverting for the observed period $T = 2.0\Gyr$:
\begin{equation}
\tauZ = f_\text{osc}\,M_\text{DM,tot}\left(\frac{2\pi}{T}\right)^2 = 7.0 \times 10^{19}\text{ J/m}^2
\end{equation}

This value, neither arbitrary nor adjusted, emerges naturally from the system's physics.

% -----------------------------------------------------------------------------
\section{MONDian Gravity: Lazy Space}
\label{sec:mond_chronology}

Beyond masses, in vast cosmic voids, spacetime becomes ``lazy''---it resists movement differently. This laziness manifests as a threshold acceleration:
\begin{equation}
a_0 = \frac{c\,H_0}{2\pi} \times \xi = 1.1 \times 10^{-10}\text{ m/s}^2
\end{equation}

The factor $\xi \simeq 1.05$ encodes the informational content of the horizon---how many quantum ``bits'' define each cell of space.

\subsection{Local Anisotropies: Mapping Tension}

Local tension variation induces variation in the Hubble ``constant'':
\begin{equation}
\frac{\delta H}{H} \simeq \frac{1}{2}\,\frac{\delta\tau}{\tauZ} \approx 10^{-4}
\end{equation}
where $\delta\tau/\tauZ$ represents the local tension contrast, estimated at about $2 \times 10^{-4}$ in the Local Supercluster vicinity. A future program capable of measuring $H_0$ directionally at 0.05\% precision over $10^\circ$ patches could reveal this cosmic tension map---regions where the membrane is tighter expand slightly faster.

% -----------------------------------------------------------------------------
\section{The Living Universe}
\label{sec:living_universe}

Our final vision: the cosmos is not an inert theater but a living organism:

\begin{table}[h]
\centering
\caption{Lifecycle of the oscillating brane universe.}
\label{tab:lifecycle}
\begin{tabular}{@{}lll@{}}
\toprule
\textbf{Phase} & \textbf{Time} & \textbf{Description} \\
\midrule
Birth & Big Bang & Maximum tension, first breath \\
Childhood & 0--1~Gyr & Relaxation, frequency tuning \\
Maturity & 1--50~Gyr & Established oscillations (\emph{we are here}) \\
Old Age & 50--100~Gyr & Progressive damping \\
Silence & $> 100$~Gyr & The strings relax, space forgets distance \\
\bottomrule
\end{tabular}
\end{table}

% -----------------------------------------------------------------------------
\section{Connection to Standard Cosmology}
\label{sec:connection_standard}

The brane framework preserves all successful predictions of \LCDM{} while adding:
\begin{enumerate}
  \item Natural explanation for dark energy timing (QCD ignition).
  \item Mechanism for MOND-like effects at large scales (Yukawa screening).
  \item Testable oscillations in cosmological observables ($w(z)$, ISW, 21\,cm).
  \item Scale-dependent structure growth reconciling DES and KiDS.
\end{enumerate}

The brane paradigm unifies inflation, dark matter, and dark energy into a single geometric framework.
